\documentclass[12pt,parskip = half, DIV=classic,listof=totoc]{scrartcl}

\usepackage[ngerman]{babel}

\usepackage[autostyle = true, german = quotes]{csquotes}
\usepackage{graphicx}
\usepackage{siunitx}
\usepackage{amsmath}
\usepackage{amssymb}
\usepackage{microtype}
\usepackage{float}
\usepackage{geometry}
\usepackage[none]{hyphenat}
\sloppy
\usepackage[labelfont={bf,sf},font={small},%
labelsep=space]{caption}
\usepackage{subcaption}
\usepackage[headsepline]{scrlayer-scrpage}
\pagestyle{scrheadings}
\automark[subsection]{section}
\ohead{\pagemark}
\chead{}
\ihead{\headmark}
\cfoot{}
\usepackage[backend=biber]{biblatex}
\addbibresource{ref-ba.bib}

\KOMAoptions{DIV = last}
\begin{document}
	\pagenumbering{gobble}
	\begin{titlepage}
		\newgeometry{top=3cm,bottom=3cm}
		\centering\bfseries\Large
		Martin-Luther-Universität\\
		Halle-Wittenberg\\
		\Large\bfseries
		Institut für Informatik
		
		\centering\vspace{3cm}\huge\bfseries
		Masterarbeit\\
		
		\vspace{3cm}\Huge
		Titel der Masterarbeit
		
		\vspace{6cm}\small\raggedright\normalfont
		\begin{tabular}{ll}
		vorgelegt von: &Max Mustermann\\
		Martrikelnummer: &123456789\\
		Studiengang: & Informatik\\
		angefertigt bei: & Prof. Dr. Beispiel\\
		Abgabedatum:  & 01.01.2026
		\end{tabular}
	\end{titlepage}
	\newgeometry{top=2.5cm,bottom=3.5cm,right=2.5cm,left=3.5cm}
	\newpage
	\tableofcontents
	\newpage
	\thispagestyle{empty}
	\section*{Eigenständigkeitserklärung}
	
	Hiermit erkläre ich, dass ich die vorliegende Arbeit eigenständig verfasst, keine anderen als die angegebenen Quellen und Hilfsmittel verwendet sowie die aus fremden Quellen direkt oder indirekt übernommenen Stellen/Gedanken als solche kenntlich gemacht habe.
	
	\vspace{4cm}
	
	\hspace{2cm} Ort, Datum \hfill Unterschrift \hspace{2cm}
	\newpage
	\pagenumbering{arabic}
	
	\section{Einleitung}
	Dies ist ein kurzer Test der Formatierung. 
	
	Die Einstellungen umfassen:
	\begin{itemize}
	    \item \texttt{scrartcl} als Dokumentenklasse (12pt, half parskip)
	    \item \texttt{polyglossia} für die deutsche Spracheinstellung
	    \item \texttt{geometry} für die Anpassung der Ränder auf 2.5cm oben/rechts, 3.5cm unten/links
	    \item \texttt{scrlayer-scrpage} für die Kopfzeilen mit Abschnittsnamen und Seitenzahlen
	\end{itemize}

	\subsection{Test Unterkapitel} Ein beispielhaftes Zitat als Test \cite{Gebauer2021_geneticDisease}.
	
	\newpage
	\printbibliography
	Hier ist ein weiterer Text, um das Layout von Unterkapiteln zu demonstrieren. Die Randeinstellungen sollten nun für diese Textseiten aktiv sein.
\end{document}